\scanpage{650}\mspage{38}
Un fait important pour la construction des
foncteurs $u^{W}_{!}$, pour \deleted{\unsure{$I$}} $u:I\longrightarrow I'$ une flèche
entre catégories ordonnées noethériennes, est
la relation
\[
 C_I\cap W(I)=(TC)_I,\tag*{(3.4.1)}
\]
d'où résulte que l'on a, pour $u$ comme dessus
\[
 u_!(C_I\cap W(I))\subset C_{I'}\cap\text{\uuline{$W(I')$}},\tag*{(3.4.2)}
\]
ingrédient essentiel de la construction du
foncteur $u^{W}_{!}$. Pour étendre cette construction
au cas d'une flèche $I$ quelconque dans Cat,
il se pose \deleted{\unsure{donc}} le problème d'étendre la
définition de $C_I$, et les formules (3.4.1) et
(3.4.2), au cas de cat. $\unsure{I|\ne}$ quelconque \unsure{dans|de} Cat.

Or on a prouvé

\uline{Proposition 3.5}. Si $\Phi\inserted{=\Psi^*}$ est une partie $Q$-\uline{close} à g.
de $\mathrm{Fl}\,M$, \deleted{\unsure{et}} et $I$ un ensemble ordonné
noethérien, $\Phi_I$ l'extension de $\Phi$ en une partie
de $\mathrm{Fl}\,M(I)$ définie par Anderson, alors on a
\[
 \Phi_I=\Psi(I)^*=v_!(\Phi(I_0)^{\sim},\tag*{\unsure{.|-}(3.5.1),}
\]
où $v:I_0\unsure{-|\to}I$ est le foncteur canonique de
la cat. discrète définie par $\mathrm{Ob}\,I$.

C'est (en des notations différentes) le
lemme 6.2.5. (p.~107) de $\underline{\overline{\mathrm{XVII}}}$, dont je ne
reprends pas ici la démonstration.

Ceci donne la clef d'une définition de $\Phi_I$
comme $\Phi\langle I\rangle$ (avec les notations des \S\S\,1,2),
% archive sheet number: 649
